\makeabstract
{
Noise: Mitigation of Uncertainty in the Precision Measurement of Long-Range Spin-Dependent Interactions
}
{
Ashley Kim (1), Nathan Clayburn (1), Yana Boneva (1), Andrew Glassford (1), Thomas Uelmen (1), Larry Hunter, Ph.D (1)
}
{
\textsuperscript{1}Department of Physics and Astronomy, Amherst College, Amherst, MA, USA
}
{
Long-range spin-dependent interactions (LRSDIs) are forces predicted by some proposed extensions to the Standard Model. We investigate the theorized effect by measuring the Larmor precession frequency of spin-polarized cesium-133 (133Cs) and mercury-199 (199Hg), composing our signal. A signal is the raw electronic data generated by optical detectors receiving electromagnetic input. Noise, in contrast, is anything that obscures the signal. While random error, uncertainty introduced by unpredictable factors, has largely been minimized, systematic error introduced by recurring discrepancies shifting the data in a consistent data is more concerning. We correct for known systematic deviations in our data and assess the remaining uncertainties associated with our correction. This summer’s work investigated various factors to further identify and minimize sources of error.
}
{
Taking data in multiple “configurations,” where parameters are altered, alerts us to systematic effects. We degauss between these different configurations to get rid of potential lingering systematic effects associated with the magnetic shields. To assess the impact of parameter variation on the precession frequencies, we measure the Hg-Cs difference field versus the slope of potential variables. We can measure the deviation of new data from historical saga data using two-sigma uncertainty analysis to identify outliers. We subtract systematic effects that we cannot physically eliminate from our final calculation of the Hg-Cs difference field.
}
{
We made several changes decreasing experimental noise: offsetting the Hg probe laser trigger to align the phase with the trough of Hg frequency best-fit error, cutting the first four measured points of observed Cs frequency from every cycle, and limiting run times to periods without high seismic noise. We shifted the phase of the electronic feedback within the 2nd harmonic generator, which improved our laser stability and potentially increased the tolerance of Hg error against construction-related noise. Additionally, we isolated the gradient coil current supply from our ground, eliminating the possibility of a ground loop.
}
{
This summer, with our primary goals being data collection and error reduction, we made significant improvements to our uncertainty and error metrics. There remain several issues, such as humidity fluctuations that correlate with the Cs field gradient, the inability to suppress third-order or higher effects in time, and lifetime decay of the cells. We plan to take data at understudied configurations to identify potential systematic errors, continue to take additional data to further improve statistical error, and run diagnostic checks on existing data.
}

\newpage
